\documentclass[11pt]{article}
\usepackage[a4paper,margin=27mm]{geometry}
\usepackage{amsmath,amssymb,amsthm,booktabs,array}
\usepackage[T1]{fontenc}
\usepackage{lmodern}
\usepackage[colorlinks=true,linkcolor=blue,citecolor=blue,urlcolor=blue]{hyperref}
\usepackage{microtype}
\newtheorem{theorem}{Theorem}
\newtheorem{proposition}[theorem]{Proposition}
\newtheorem{corollary}[theorem]{Corollary}
\newcommand{\R}{\mathbb R}
\newcommand{\G}{\mathcal G}
\title{Resolving Question III.5 of Stepanov--Herrmann--Guhr\\
(arXiv:2104.08966v1, p.~7):\\
the top-eigenvector weight gap occurs exactly above $\sigma=c$}
\author{Haobo Ma\\\small ChronoAI Pte Ltd, Singapore
\and Wenlin Zhang\\\small National University of Singapore\\
\small ChronoAI Pte Ltd, Singapore}
\date{3 November 2026}
\begin{document}
\maketitle
\begin{abstract}
We resolve Question III.5 of Stepanov, Herrmann and Guhr
(arXiv:2104.08966v1, p.~7): on their remaining domain
$A\setminus(A_1\cup A_2)$, correlation matrices with prescribed
off-diagonal mean $c$ and standard deviation $\sigma$ and a nonmaximal
top-eigenvector weight exist in every sufficiently large dimension
exactly when $\sigma>c$.
A centered spectral-moment inequality excludes $\sigma=c$ in every
dimension. A rank-two perturbation of a scalar matrix gives the required
matrices for $\sigma>c$. The same construction gives an exact eventual
classification on the entire legal half-disc, including zero mean and
the circular boundary. An eigenspace criterion makes the requirement
of a weight gap in every eigenbasis explicit.
\end{abstract}

\section{The question and its conventions}
We answer Question III.5 of \cite[p.~7]{SHG} with the basis-independent
meaning of the weight gap in their Definition VI.1 (p.~12).
The question is answered by Theorem~\ref{thm:question}; the boundary
corollaries below complete the eventual classification.
A correlation matrix is a real symmetric positive semidefinite matrix
$C$ with $C_{ii}=1$. For $n\ge2$, put
\begin{equation}\label{eq:moments}
 c(C)=\frac{2}{n(n-1)}\sum_{i>j}C_{ij},\qquad
 \sigma(C)^2=\frac{2}{n(n-1)}\sum_{i>j}C_{ij}^2-c(C)^2,
\end{equation}
where $\sigma(C)\ge0$. If $Cv_j=\lambda_jv_j$ is an orthonormal
eigenbasis, ordered by $\lambda_1\ge\cdots\ge\lambda_n$, set
$\delta_n=(1,\ldots,1)/\sqrt n$ and $w_j=\langle v_j,\delta_n\rangle^2$.
Write $\G(C)$ if, for every such basis and every index of a maximal
eigenvalue, another index has a strictly larger weight.
This includes all choices inside repeated eigenspaces.

The source defines
\begin{align*}
 A&=\{(c,\sigma):c>0,\ \sigma>0,\ c^2+\sigma^2<1\},\\
 A_1&=\{(c,\sigma)\in A:c\ge\tfrac12\ \text{or}\ \sigma<c\},\\
 A_2&=\{(c,\sigma)\in A:\sigma>\sqrt3c\ \text{or}\ c<\sigma<1-c\}.
\end{align*}
Let $D=A\setminus(A_1\cup A_2)$. Its points satisfy
$0<c<1/2$, $\sigma\ge c$, and either $\sigma=c$ or
$1-c\le\sigma\le\sqrt3c$.

\begin{samepage}
\begin{theorem}\label{thm:question}
For each $(c,\sigma)\in D$, the following are equivalent:
\begin{enumerate}
\item There is an $N$ such that every integer $n\ge N$ admits an
$n\times n$ correlation matrix $C$ with
$c(C)=c$, $\sigma(C)=\sigma$, and $\G(C)$.
\item $\sigma>c$.
\end{enumerate}
If $\sigma=c$, no such matrix exists in any dimension $n\ge2$.
\end{theorem}
\end{samepage}

The characteristic identities in \cite[Lemma II.5]{SHG} supply the
spectral moments used below. The construction has two constant
within-group correlations and one between-group correlation. Positive
definiteness of this block-matrix class is also studied in \cite{RD};
here a sum of positive rank-one matrices makes positivity explicit.

\section{Obstructions from two spectral moments}
\begin{proposition}\label{prop:obstruction}
Let $C$ have size $n\ge3$ and $c=c(C)>0$, $\sigma=\sigma(C)$.
If a top eigenvector in an orthonormal eigenbasis has weight at most
$1/2$, then
\begin{equation}\label{eq:necessary}
 \sigma^2\ge\frac{n}{n-2}c^2,\qquad
 c\le\frac{n-2}{2(n-1)}.
\end{equation}
Consequently $\G(C)$ is impossible if $\sigma\le c$ or $c\ge1/2$.
\end{proposition}
\begin{proof}
Put $\mu_j=\lambda_j-1$. The trace, squared Frobenius norm, and
Rayleigh quotient give
\begin{equation}\label{eq:spectral}
 \sum_j\mu_j=0,\quad
 \sum_j\mu_j^2=n(n-1)(c^2+\sigma^2),\quad
 \sum_j\mu_jw_j=(n-1)c,\quad \sum_jw_j=1.
\end{equation}
For example, $\operatorname{tr}(C^2)=n+n(n-1)(c^2+\sigma^2)$,
and $\delta_n^TC\delta_n=1+(n-1)c$, proving these identities.
Take a top index as index 1 and the largest remaining eigenvalue as
index 2. If $w_1\le1/2$, then
\[
 (n-1)c\le\mu_1w_1+\mu_2(1-w_1)
 \le\frac{\mu_1+\mu_2}{2}.
\]
Set $S=\mu_1+\mu_2>0$. Since $\sum_{j\ge3}\mu_j=-S$,
Cauchy--Schwarz gives
\[
 \sum_j\mu_j^2\ge\frac{S^2}{2}+\frac{S^2}{n-2}
 =\frac{nS^2}{2(n-2)}
 \ge\frac{2n(n-1)^2c^2}{n-2}.
\]
Substitution into \eqref{eq:spectral} proves the first inequality
in \eqref{eq:necessary}. Independently, nonnegativity and trace $n$
give
\[
 1+(n-1)c=\sum_j\lambda_jw_j
 \le\tfrac12(\lambda_1+\lambda_2)\le n/2,
\]
which proves the second. A strict weight gap forces $w_1<1/2$:
if $w_j>w_1$ for some $j\ne1$, then $2w_1<w_1+w_j\le1$.
For $n=2$, there is a single off-diagonal coefficient, so
$\sigma(C)=0$; if $c>0$, its top eigenvector is $\delta_2$ and has
weight 1. This also excludes the relevant two-dimensional cases.
\end{proof}

\begin{corollary}\label{cor:majority}
If $c>0$ and $\sigma\le c$, the largest eigenvalue is simple and its
weight is greater than $1/2$ in every dimension. Its acute angle
$\theta_1=\arccos\sqrt{w_1}$ with $\delta_n$ satisfies
$\theta_1\le\arctan(\sigma/c)$.
\end{corollary}
\begin{proof}
For $n\ge3$, Proposition~\ref{prop:obstruction} excludes $w_1\le1/2$.
Any unit top eigenvector extends to an orthonormal eigenbasis. If the
top eigenspace had dimension at least two, it would contain a unit
vector orthogonal to $\delta_n$, contradicting $w_1>1/2$.
The case $n=2$ follows from the last paragraph of that proposition.
To obtain the angle bound, center the two vectors
$\mu=(\lambda_j-1)_j$ and $w=(w_j)_j$ in \eqref{eq:spectral}.
Cauchy--Schwarz gives
\[
 (n-1)^2c^2\le n(n-1)(c^2+\sigma^2)
 \left(\sum_jw_j^2-\frac1n\right).
\]
Thus $\sum_jw_j^2\ge \frac1n+\frac{n-1}{n}
\frac{c^2}{c^2+\sigma^2}\ge\frac{c^2}{c^2+\sigma^2}$,
which is also \cite[Eq.~(IV.17)]{SHG}.
Since $w_1>1/2$, $w_1=\max_jw_j\ge\sum_jw_j^2$.
Taking square roots and acute angles proves the assertion.
This extends the sufficient angular regime of Remark III.7
in \cite[p.~8]{SHG} to $\sigma\le c$.
\end{proof}

In particular, any positive-mean matrix with $\G$ must have
\[
 n\ge\max\left\{\frac{2\sigma^2}{\sigma^2-c^2},
                   \frac{2(1-c)}{1-2c}\right\}.
\]
Both denominators are positive by Proposition~\ref{prop:obstruction};
the bounds follow by rearranging \eqref{eq:necessary}.

\section{An exact rank-two construction}
\begin{proposition}\label{prop:construction}
If $0<c<1/2$, $\sigma>c$, and $c^2+\sigma^2<1$, then in every
sufficiently large dimension there is a positive definite correlation
matrix with characteristic $(c,\sigma)$ satisfying $\G$.
\end{proposition}
\begin{proof}
Define
\[
 L=\frac{\sigma^2-(1-c)^2}{\sigma^2+(1-c)^2},\qquad
 U=\frac{\sigma^2-c^2}{\sigma^2+c^2},\qquad
 x=\frac{\max(0,L)+\min(c,U)}2.
\]
The interval containing $x$ is nonempty. Indeed $U>0$,
$L<c$ follows from $(1-c)(1-c^2-\sigma^2)>0$, and
$L<U$ follows from $2\sigma^2(1-2c)>0$ after multiplication by
the positive denominators. Thus
\begin{equation}\label{eq:interval}
 \max(0,L)<x<\min(c,U)<1/2.
\end{equation}
Set $a=\sigma/\sqrt{1-x^2}$, $b=c-ax$, and $d=1-a-b$.
Then $a,b,d>0$ and $a+b>2c$.
For $b>0$, observe that
$(c^2+\sigma^2)x^2<c^2$, since $x<c$ and $c^2+\sigma^2<1$.
The other two inequalities follow by squaring positive quantities:
\[
 d>0\ \Longleftrightarrow\
 \sigma^2\frac{1-x}{1+x}<(1-c)^2\ \Longleftrightarrow\ x>L,
\]
\[
 a+b>2c\ \Longleftrightarrow\
 \sigma^2\frac{1-x}{1+x}>c^2\ \Longleftrightarrow\ x<U.
\]

For each $n\ge2$, choose a nearest integer $k_n$ to
$n(1+\sqrt{x})/2$, and let $s\in\{1,-1\}^n$ have $k_n$ plus
signs. Write $e=(1,\ldots,1)^T$ and
\begin{align*}
 t_n&=\frac{2k_n-n}{n},&
 z_n&=\frac{(2k_n-n)^2-n}{n(n-1)},\\
 a_n&=\frac{\sigma}{\sqrt{1-z_n^2}},&
 b_n&=c-a_nz_n,& d_n&=1-a_n-b_n.
\end{align*}
The rounding error is at most $1/2$, so $t_n\to\sqrt{x}$ and
$z_n=(nt_n^2-1)/(n-1)\to x$.
Continuity and the strict inequalities above give, for all sufficiently
large $n$, $0<k_n<n$, $t_n>0$, $|z_n|<1$, $a_n,b_n,d_n>0$, and
\begin{equation}\label{eq:finitegap}
 2(n-1)c<(n-2)(a_n+b_n).
\end{equation}
All formulas below are used only for these dimensions. Put
\begin{equation}\label{eq:matrix}
 C_n=d_n I+a_nss^T+b_nee^T.
\end{equation}
It is positive definite, with diagonal $d_n+a_n+b_n=1$.
The average of $s_is_j$ over $i>j$ equals $z_n$, and
$(s_is_j)^2=1$. Hence its mean is $a_nz_n+b_n=c$ and its variance is
\[
 (a_n^2+b_n^2+2a_nb_nz_n)-(a_nz_n+b_n)^2
 =a_n^2(1-z_n^2)=\sigma^2.
\]

It remains to establish the weight gap, including all eigenbasis choices.
Let $B_n=C_n-d_nI$. Since $a_n,b_n>0$ and $s,e$ are independent,
$B_n$ has two positive eigenvalues $\gamma_+>\gamma_->0$.
For distinctness, in the orthonormal basis starting with $e/\sqrt n$
on $\operatorname{span}(e,s)$, the off-diagonal entry is
$na_nt_n\sqrt{1-t_n^2}>0$. All other eigenvalues of $B_n$ are zero;
their eigenvectors are orthogonal to $e$ and have weight zero.
The two positive eigendirections therefore carry weights $p_+,p_-$
with $p_++p_-=1$. Their weighted mean and their eigenvalue sum are
\[
 r:=\delta_n^TB_n\delta_n
 =a_n\frac{(2k_n-n)^2}{n}+b_nn
 =(n-1)c+a_n+b_n,\qquad
 T:=\gamma_++\gamma_-=n(a_n+b_n).
\]
By \eqref{eq:finitegap}, $r<T/2$. The identity
\[
 r-\frac T2=(\gamma_+-\gamma_-)(p_+-\tfrac12)
\]
implies $p_+<1/2<p_-$. The two eigendirections are simple and
all remaining weights vanish, proving $\G(C_n)$ for every eigenbasis.
\end{proof}

\begin{proof}[Proof of Theorem~\ref{thm:question}]
On $D$, $\sigma\ge c$ and $0<c<1/2$.
Proposition~\ref{prop:obstruction} excludes $\sigma=c$ in every
dimension, while Proposition~\ref{prop:construction} supplies a matrix
in every sufficiently large dimension when $\sigma>c$.
\end{proof}

\begin{corollary}\label{cor:all}
For any $(c,\sigma)\in A$, eventual existence with $\G$ holds
exactly when $c<1/2$ and $\sigma>c$.
For every compact subset of this open region, one common dimension
threshold works for all its points.
\end{corollary}
\begin{proof}
The two propositions prove the pointwise assertion.
For compact uniformity, $x$, $a$, $b$, $d$, and $a+b-2c$ are continuous
functions of $(c,\sigma)$ and their required positive margins have
positive minima on the compact set. Also $x$ stays in a compact
subinterval of $(0,1/2)$. The bound
$|t_n-\sqrt{x}|\le1/n$ makes $z_n\to x$ uniform there. The
coefficients and the strict inequality \eqref{eq:finitegap} therefore
hold uniformly beyond one threshold.
\end{proof}

\section{Eigenspaces and boundary characteristics}
The universal quantifier in Definition VI.1 matters when eigenvalues
are repeated. The following criterion applies to every real symmetric
matrix, with weights relative to any fixed unit vector $u$.
\begin{proposition}\label{prop:projectors}
Let $E_0$ be the top eigenspace, and let $E_1,\ldots,E_r$ be the
eigenspaces for the distinct smaller eigenvalues. Put
$q_j=\|P_{E_j}u\|^2$ and $m_j=\dim E_j$.
The universal weight gap holds if and only if
\begin{equation}\label{eq:projectors}
 q_0<\max_{1\le j\le r}\frac{q_j}{m_j},
\end{equation}
where the maximum of an empty family is zero.
\end{proposition}
\begin{proof}
In any orthonormal basis of $E_j$, the squared components of
$P_{E_j}u$ sum to $q_j$. Hence one is at least $q_j/m_j$,
whereas each weight in $E_0$ is at most $q_0$. This proves
sufficiency. Conversely, choose a basis of $E_0$ concentrating its
projection in one vector, whose weight is $q_0$. In each lower
$E_j$, choose an orthonormal basis in which all projection
components have absolute value $\sqrt{q_j/m_j}$. Such a basis
exists: an orthogonal map sends the projection's unit direction to
$(1,\ldots,1)/\sqrt{m_j}$; if $q_j=0$ every basis works.
When \eqref{eq:projectors} fails, no weight in this combined basis
exceeds the chosen top weight. The universal gap therefore fails.
The same reasoning covers the scalar-matrix case $r=0$.
\end{proof}

\begin{corollary}\label{cor:circle}
For $n\ge2$, a correlation matrix with $c^2+\sigma^2=1$ satisfies
$\G$ if and only if $c<0$. In particular, none on this boundary
with $c\ge0$ satisfies the universal gap.
\end{corollary}
\begin{proof}
All off-diagonal entries have absolute value at most one. Equality of
their average squared value to one forces all to have absolute value
one. In a unit-vector Gram representation, equality in
Cauchy--Schwarz makes all vectors signed copies of a single vector.
Thus $C=ss^T$ with $s_i\in\{-1,1\}$, as in
\cite[Lemma VI.11 and Example II.4]{SHG}.
Here $E_0=\R s$, $m_1=n-1$, and
$q_0=(\sum_i s_i)^2/n^2=(1+(n-1)c)/n$.
Criterion~\eqref{eq:projectors} is $q_0<(1-q_0)/(n-1)$,
equivalently $c<0$.
\end{proof}

\begin{corollary}\label{cor:zero}
For every fixed $0<\sigma<1$, positive definite correlation matrices
with $c=0$, standard deviation $\sigma$, and $\G$ exist in every
sufficiently large dimension.
\end{corollary}
\begin{proof}
In \eqref{eq:matrix}, take $k=\lceil n/2\rceil$ plus signs and set
\[
 z=\begin{cases}-1/(n-1),&n\text{ even},\\-1/n,&n\text{ odd},\end{cases}
 \qquad a=\frac{\sigma}{\sqrt{1-z^2}},\quad b=-az,\quad d=1-a-b.
\]
For sufficiently large $n$, $a,b,d>0$ because $a\to\sigma$,
$b\to0$ and $d\to1-\sigma>0$. The moment calculation of
Proposition~\ref{prop:construction} gives exactly $c=0$ and
standard deviation $\sigma$. The rank-two trace and Rayleigh
quotient there are $T=n(a+b)$ and $r=a+b$, so $2r<T$ for $n>2$.
For even $n$, $e\perp s$ and the active eigenvalues $na,nb$ are
distinct because $a>b$. For odd $n$, their off-diagonal coupling
is positive as in that proposition. In both cases the active
directions are simple, their weights sum to one, and
$r-T/2=(\gamma_+-\gamma_-)(p_+-1/2)<0$ proves $\G$.
The zero-mean examples in \cite[Section VII, p.~18]{SHG} thus have
an exact prescribed-variance construction in every large dimension.
\end{proof}

\begin{corollary}\label{cor:legal}
Fix $(c,\sigma)$ in the legal half-disc
$\{(c,\sigma):\sigma\ge0,\ c^2+\sigma^2\le1\}$.
Eventual existence with $\G$ holds exactly when
\[
 0\le c<1/2,\qquad \sigma>c,\qquad c^2+\sigma^2<1.
\]
\end{corollary}
\begin{proof}
For $c>0$ in the interior with $\sigma>0$, use
Corollary~\ref{cor:all}; for $c=0$ use Corollary~\ref{cor:zero}.
If $c<0$, positive semidefiniteness gives
$0\le\delta_n^TC\delta_n=1+(n-1)c$, which fails for large $n$.
For $c\ge0$ on the circle use Corollary~\ref{cor:circle}.
If $\sigma=0$ and $c>0$, then
$C=(1-c)I+cee^T$ has top weight one; if $c=\sigma=0$, $C=I$
has an eigenbasis containing $\delta_n$, so $\G$ fails.
These cases exhaust the legal half-disc.
\end{proof}

The same centered argument yields a spectral obstruction beyond
correlation matrices, addressing the wider-matrix direction mentioned
in \cite[Remark IV.7 and Section VII]{SHG}.
\begin{corollary}\label{cor:symmetric}
Let $H$ be real symmetric of size $n\ge3$, let $u$ be a unit vector,
$\tau=\operatorname{tr}H/n$, and $r=u^THu-\tau>0$.
If a top eigenvector has weight at most $1/2$ relative to $u$, then
\[
 \operatorname{tr}\bigl((H-\tau I)^2\bigr)\ge\frac{2n}{n-2}r^2.
\]
If also $H$ is positive semidefinite, then
$u^THu\le\operatorname{tr}H/2$.
\end{corollary}
\begin{proof}
Center the eigenvalues at $\tau$. They sum to zero, and their
weighted sum is $r$. The pair estimate and Cauchy--Schwarz in
Proposition~\ref{prop:obstruction} give the first assertion with
$(n-1)c$ replaced by $r$. Nonnegative eigenvalues give
$u^THu\le(\lambda_1+\lambda_2)/2\le\operatorname{tr}H/2$.
\end{proof}

\section{A rational example and coverage}
Take $n=16$, $k=12$, $a=2/3$, $b=1/4$, $d=1/12$ in
\eqref{eq:matrix}. Then $z=1/5$,
\[
 c=\frac{23}{60},\qquad \sigma^2=\frac{32}{75},\qquad
 \lambda_1=\frac{145}{12},\quad\lambda_2=\frac{11}{4},\quad
 \lambda_3=\cdots=\lambda_{16}=\frac1{12}.
\]
This point lies in $D$ above the diagonal. The top two eigenvectors
are proportional to $e+4s$ and $3e-2s$, respectively, and their weights
are exactly $3/7$ and $4/7$.
The accompanying \texttt{verify.py} checks this example, exact moment
and rank-two certificates in dimensions $16,17,24,100$, and centered
two-point energy inequalities for integer controls in dimensions
$3$ through $30$. It also checks exact zero-mean and circular-boundary
controls in dimensions $4$ through $40$.
These finite checks illustrate the formulas;
the propositions prove the quantified statements.

\begin{center}\small
\begin{tabular}{@{}>{\raggedright\arraybackslash}p{0.34\linewidth}>{\raggedright\arraybackslash}p{0.59\linewidth}@{}}
\toprule
Source item in \cite{SHG} & Coverage category and result \\
\midrule
Question III.5, p.~7 & Proved: Theorem~\ref{thm:question}, exact answer on $D$.\\
Theorems III.3--III.4, pp.~6--7 & Prior proofs cited for $A_1,A_2$;
proved completion and compact uniformity in Corollary~\ref{cor:all}.\\
Remark III.7, p.~8 & Proved: the top-direction angle conclusion extends
to $\sigma\le c$, Corollary~\ref{cor:majority}.\\
Definition VI.1, p.~12 & Proved: exact eigenspace criterion,
Proposition~\ref{prop:projectors}, for every symmetric matrix.\\
Example II.4, p.~4; Lemma VI.11, p.~15 & Prior rank-one description cited;
proved universal gap exactly for $c<0$, Corollary~\ref{cor:circle}.\\
Section VII, p.~18: $c=0$ and $c<0$ & Proved: exact zero-mean witnesses
in every large dimension and full legal-half-disc classification,
Corollaries~\ref{cor:zero}--\ref{cor:legal}.\\
Remark IV.7, p.~10; Section VII, p.~18 & Prior largest-eigenvalue estimate
cited; proved centered weight obstruction for symmetric matrices,
Corollary~\ref{cor:symmetric}.\\
Section VII, p.~18: skewness/kurtosis & Not covered: higher-moment
refinements of the source's spectral estimates.\\
\bottomrule
\end{tabular}
\end{center}

\begingroup\raggedright
\paragraph{Formal scope.}
In the pinned
\href{https://github.com/the-omega-institute/trureturing/blob/8c3421e2977f07cb0b5d9b30b7ede138d9951cb5/D5/S3/Estimation/CorrelationTopEigenvectorWeight.lean}{\texttt{CorrelationTopEigenvectorWeight.lean}},
the declaration in namespace \texttt{D5.S3.Estimation},
\texttt{CorrelationTopEigenvectorWeight.result},
proves \texttt{claim}: for each point of \texttt{domainD}, the eventual
existence condition of Theorem~\ref{thm:question} is equivalent to $c<\sigma$.
Its \texttt{IsCorrelation}, \texttt{meanCorr}, \texttt{stdCorr} and
\texttt{SatisfiesWGap} definitions match \eqref{eq:moments} and $\G$.
The explicit finite-dimensional bounds, all-$A$ and legal-half-disc
classifications, compact uniformity, majority/angle bound, eigenspace
criterion, boundary corollaries, general symmetric-matrix obstruction,
and rational examples have the written proofs above; their additional
scope is distinct from that declaration. The checker covers the finite
identities and parameter ranges specified alongside the example.\par
\endgroup

\begin{thebibliography}{9}
\bibitem{SHG}
Y.~Stepanov, H.~Herrmann and T.~Guhr,
\emph{Generic features in the spectral decomposition of correlation matrices},
J. Math. Phys. \textbf{62} (2021), 083505,
\href{https://doi.org/10.1063/5.0054438}{doi:10.1063/5.0054438}.
The question, numbering and page locators here refer to
\href{https://arxiv.org/abs/2104.08966v1}{arXiv:2104.08966v1}.
\bibitem{RD}
O.~Roustant and Y.~Deville,
\emph{On the validity of parametric block correlation matrices with constant
within and between group correlations},
\href{https://arxiv.org/abs/1705.09793v1}{arXiv:1705.09793v1} (2017).
\end{thebibliography}
\noindent\small Licensed under Creative Commons Attribution 4.0 International
(CC BY 4.0; \texttt{cc-by-4.0}).
\end{document}
